\renewcommand{\Anno}{2009}

\renewcommand{\NumeroProblema}{3}

\renewcommand{\Testo}{ 
Marco ha costruito 21 tessere simili a quelle del domino.  Ogni
tessera \`{e} un rettangolo suddiviso in due quadratini, ed ogni
quadratino contiene un numero di puntini colorati compreso tra 0 e 5
(estremi inclusi).  Non ci sono due tessere uguali (come numero di
puntini nei due quadratini), ed ogni coppia di numeri
rappresenta il numero dei puntini di esattamente una tessera.

Ora Marco vorrebbe disporre alcune di queste tessere in una fila
orizzontale, facendo in modo che 
\begin{itemize}
    \item due tessere consecutive abbiano lo stesso numero di puntini
    sui quadratini che si toccano,

    \item i due quadratini all'estrema destra ed all'estrema sinistra
    della fila abbiano lo stesso numero di puntini.
\end{itemize}

Quello che segue \`{e} un esempio di una fila orizzontale costituita
da 7 tessere che rispetta entrambe le condizioni (i numeri
rappresentano i puntini presenti in quel quadratino).
$$
\domino{0}{1}-\domino{1}{1}-\domino{1}{3}-\domino{3}{5}-\domino{5}{4}-
\domino{4}{3}-\domino{3}{0}$$
\begin{enumerate}
    \item Determinare quante tessere al massimo possono far parte di
    una fila che rispetta entrambe le condizioni.

    \item Rispondere alla stessa domanda nel caso in cui le tessere,
    invece di essere tutte quelle possibili con numero di puntini da 0
    a 5, siano tutte quelle possibili con numero di puntini da 0 a
    2009.
\end{enumerate}

{\small Nota bene: c'\`{e} un'unica tessera, per esempio, che ha 3 e 4
puntini; tale tessera pu\`{o} essere utilizzata sia in versione
$\domino{3}{4}$, sia in versione $\domino{4}{3}$. 
}}

\renewcommand{\Risposta}{
\begin{enumerate}
	\renewcommand{\labelenumi}{(\alph{enumi})} 
	\item  18

	\item  $2\cdot 1005^{2}$
\end{enumerate}
}

\renewcommand{\Aiutino}{
\begin{enumerate}
	\renewcommand{\labelenumi}{(\alph{enumi})} 
	\item Quanti sono complessivamente i quadratini che contengono
	ciascuno dei numeri da 0 a 5? In ogni fila con le propriet\`{a}
	richieste ogni quadratino compare un numero pari di volte
	(perch\'{e}?). Ne segue che un po' di quadratini devono restare
	esclusi. In questo modo si dimostra che al massimo 18 tessere
	possono far parte della fila. Non bisogna poi dimenticare di
	fornire un esempio di fila con 18 tessere.

	\item Per dimostrare che un certo numero di tessere devono essere
	escluse basta estendere il ragionamento del caso precedente.  Per
	fornire una costruzione si pu\`{o} procedere per induzione,
	mostrando come una una fila massinale per un certo numero dispari
	$2n+1$ pu\`{o} essere allungata per produrre una fila massimale
	per il dispari $2n+3$.

\end{enumerate}
}

\renewcommand{\SoluzioneUfficiale}{\DomandaA La pi\`{u} lunga fila
possibile che soddisfi entrambe le condizioni contiene 18 tessere.

Fissiamo un qualunque intero $k$ da 0 a 5.  Il numero totale di
quadratini con $k$ punti colorati nelle varie tessere \`{e} 7 (non
bisogna dimenticare la tessera con i 2 quadratini uguali, la quale
conta doppio).  Inoltre, in una qualunque fila orizzontale che
rispetti i vincoli, il numero totale di quadratini con $k$ punti
colorati \`{e} pari.  Infatti per costruzione i quadratini compaiono
sempre ``accoppiati'', sia quelli lungo la fila, sia quelli ai 2
bordi.

Ne segue che almeno uno dei 7 quadratini con $k$ punti colorati deve
non comparire nella fila.  Questo ragionamento vale per tutti i 6
possibili valori di $k$.  Pertanto almeno 6 quadratini, uno per
numero, dovranno mancare dalla fila, il che vuol dire che almeno 3
tessere resteranno inutilizzate.

Questo ragionamento mostra che la fila pi\`{u} lunga potr\`{a}
contenere al pi\`{u} $21-3=18$ tessere.  Per concludere dobbiamo ora
esibire un esempio di fila lunga 18.  Scartate le tessere
$$\domino{0}{1}\quad \domino{2}{3}\quad \domino{4}{5}$$ 
(si noti che ogni numero compare una ed una sola volta), le restanti
possono essere disposte ad esempio nel seguente modo:
$$\domino{0}{5}-\domino{5}{5}-\domino{5}{1}-\domino{1}{4}-\domino{4}{4}
-\domino{4}{2}-\domino{2}{5}-\domino{5}{3}-\domino{3}{4}-
\domino{4}{0}-\mbox{}$$
$$\mbox{}-\domino{0}{0}-\domino{0}{2}-\domino{2}{2}-\domino{2}{1}-
\domino{1}{1}-\domino{1}{3}-\domino{3}{3}-\domino{3}{0}$$

\DomandaB Mostriamo pi\`{u} in generale che, se le tessere contengono 
i numeri da 0 a $2n+1$ (dispari!), la pi\`{u} lunga fila compatibile
con i vincoli contiene $2(n+1)^{2}$ tessere. Nel caso $n=2$ si
ritrovano le 18 tessere precedenti. Nel caso $n=1004$ (che corrisponde
ad avere i numeri da 1 fino a 2009) si trovano $2\cdot 1005^{2}$
tessere.

Ragionando come nel punto (a) si ha che ogni $k$ compare in $(2n+3)$
quadratini, quindi, trattandosi di un numero dispari, almeno uno di
tali quadratini non potr\`{a} comparire nella fila.  Poich\'{e} il
ragionamento vale per tutti i $(2n+2)$ valori possibili di $k$,
abbiamo che dalla fila dovranno mancare almeno $(2n+2)$ quadratini,
dunque almeno $(n+1)$ tessere.

Non resta che contare le tessere totali.  Ce ne sono $(2n+2)$ con i 2
quadratini uguali, pi\`{u} $(2n+2)(2n+1)/2$ con i 2 quadratini diversi
(abbiamo diviso per 2 perch\'{e} ogni coppia di numeri distinti
compare una sola volta). Dunque le tessere in totale sono
$$(2n+2)-\frac{(2n+2)(2n+1)}{2}=2n^{2}+5n+3.$$

Togliendo le $(n+1)$ tessere che di sicuro non possono far parte della
fila, restano al pi\`{u} 
$$2n^{2}+5n+3-(n+1)=2(n+1)^{2}$$ 
tessere nella fila.

Per concludere bisogna ora esibire un esempio di fila con questa
lunghezza. Per far questo procediamo per induzione. Per $n=1$ abbiamo 
la fila
$$\domino{0}{0}-\domino{0}{2}-\domino{2}{2}-\domino{2}{1}-
\domino{1}{1}-\domino{1}{3}-\domino{3}{3}-\domino{3}{0}$$

Supponiamo ora di avere una fila corrispondente ad un certo $n$
(contenente quindi $2(n+1)^{2}$ tessere con numeri da $1$ a $2n+1$), e
mostriamo come costruire una fila corrispondente ad $(n+1)$
(contenente quindi $2(n+2)^{2}$ tessere con numeri da $1$ a $2n+3$).
A tal fine supponiamo, senza perdita di generalit\`{a}, che la fila da
cui partiamo inizi e termini con uno 0, ed aggiungiamo
\emph{all'inizio} tutte le tessere con i numeri $(2n+2)$ e $(2n+3)$
meno una.

Iniziamo dunque la fila con il seguente schema
$$\domino{0}{2n+3}-\domino{2n+3}{1}-\domino{1}{2n+2}-\domino{2n+2}{2}-\mbox{}$$
$$\mbox{}-\domino{2}{2n+3}-\domino{2n+3}{3}-
\domino{3}{2n+2}-\domino{2n+2}{4}-\mbox{\ldots}$$
e ripetiamo lo schema di 4 in 4 fino a concludere con
$$\mbox{}-\domino{2n}{2n+3}-\domino{2n+3}{2n+1}-
\domino{2n+1}{2n+2}-\domino{2n+2}{0}-\mbox{}$$

Da qui in poi procediamo come nel caso $n$. In questo modo abbiamo
piazzato tutte le tessere con i numeri $(2n+2)$ e $(2n+3)$,
escluse le 3 tessere
$$\domino{2n+2}{2n+3}\quad\domino{2n+2}{2n+2}\quad\domino{2n+3}{2n+3}$$

La prima \`{e} quella che decidiamo di ``sacrificare'', le restanti 2
si possono agevolmente inserire in molti modi nella prima parte della
fila.

Per fare un esempio di come funziona la costruzione in un caso
concreto, passando da $n=1$ ad $n=2$ si ottiene la fila di 18 tessere 
che abbiamo esibito al termine della soluzione della parte~(a).
}

\renewcommand{\SoluzioniAlternative}{\TitoloSoluzioneAlternativaA
Risolviamo direttamente il caso dispari generale, cio\`{e} quello in
cui nei quadratini sono rappresentati tutti i possibili numeri da 0 a 
$2n+1$. Osserviamo intanto che le tessere con 2 quadratini uguali non 
sono influenti, in quanto si possono facilmente aggiungere o togliere 
da una fila senza alterarne le propriet\`{a}. Il problema \`{e}
pertanto equivalente a trovare la fila massimale che utilizza solo
tessere composte da quadratini con numeri di punti diversi.

Interpretiamo ora il problema in termini di teoria dei grafi.
Costruiamo un grafo in cui i vertici sono i numeri da 0 a $(2n+1)$, e 
2 vertici sono collegati da un lato se e solo se c'\`{e} una tessera
che ha quei 2 numeri rappresentati. Avendo escluso le tessere con 2
quadratini uguali, sappiamo che ogni lato connette 2 vertici distinti,
come nella definizione standard di grafo. Sapendo poi che abbiamo
tutte le tessere possibili, abbiamo che nel grafo ci sono tutti i
lati possibili. Abbiamo quindi quello che si chiama un \emph{grafo
completo}, con $(2n+2)$ vertici e $\binom{2n+2}{2}$ lati.

Una fila di tessere con le 2 propriet\`{a} proposte corrisponde nel
grafo ad un cammino che parte da un vertice e torna a quello stesso
vertice dopo aver percorso un certo numero di lati del grafo (ogni lato 
viene percorso al pi\`{u} una volta). 

Ora \`{e} ben noto che in un grafo un cammino chiuso passa sempre per 
un numero \emph{pari} di lati uscenti da ogni vertice (questo perch\'{e} ogni
volta che un cammino arriva in un vertice mediante un certo lato, poi 
riparte seguendo un altro lato). Nel nostro caso in ogni vertice
concorrono $(2n+1)$ lati, dunque almeno uno di essi non far\`{a} parte
del cammino. Ripetendo lo stesso ragionamento per tutti i $(2n+2)$
vertici, abbiamo l'esclusione di almeno $(n+1)$ lati (abbiamo diviso
per 2 perch\'{e} ogni lato escluso ``accontenta'' 2 vertici).

Resta ora da fare la seconda parte, cio\`{e} mostrare che esiste un
cammino chiuso che esclude soltanto $(n+1)$ lati.  Per far questo
osserviamo intanto che \`{e} possibile eliminare $(n+1)$ lati in modo
da coinvolgere una ed una sola volta tutti i vertici: basta ad esempio
eliminare il lato da 0 a 1, quello da 2 a 3, e cos\`{i} via.  Al
termine dell'operazione ci resta un grafo in cui \emph{da ogni vertice
parte un numero pari di lati}.  Un ben noto teorema (non banale) della
teoria dei grafi ci dice che questa \`{e} la condizione necessaria e
sufficiente per l'esistenza di un cammino chiuso che passa una ed una
sola volta per tutti i lati del grafo (circuito hamiltoniano). Questo 
teorema generale ci permette quindi di concludere la dimostrazione
sostanzialmente gratis!
}

\renewcommand{\Osservazioni}{\TitoloOsservazione
\`{E} interessante esaminare il caso in cui le tessere contengono i
numeri da 1 a $2n$ (pari!). In tal caso ogni $k$ compare in $(2n+2)$
quadratini, che \`{e} un numero pari. Cade dunque l'ostruzione alla
presenza di tutti i quadratini nella fila che c'era nel caso dispari.

In effetti nel caso pari esiste sempre una fila che rispetta le 2
condizioni e contiene tutte le possibili tessere.  La dimostrazione
dell'esistenza di tale fila si pu\`{o} fare per induzione riadattando,
con minime modifiche, la stessa costruzione proposta nel caso dispari
(in questo caso non serve nemmeno ``sacrificare'' una tessera nel
passo induttivo).

In termini di teoria dei grafi, in questo caso abbiamo il grafo
completo su $(2n+1)$ vertici.  In ogni vertice concorrono $2n$ lati,
dunque un numero pari, e pertanto esiste un circuito hamiltoniano.  }

\renewcommand{\Commenti}{\TitoloCommenti \`{E} importante notare che
questo \`{e} un classico esempio di problema in cui ci sono 2
disuguaglianze da dimostrare.  Nella parte (a), ad esempio, bisogna
sia dimostrare che la fila pi\`{u} lunga contiene \emph{al massimo} 18
tessere (ad esempio con l'argomento di parit\`{a} che abbiamo
proposto), sia dimostrare che la fila pi\`{u} lunga contiene
\emph{almeno} 18 tessere (per questo basta esibire un esempio).

Entrambe queste parti sono necessarie per la risoluzione completa
della parte (a). Stesso discorso vale per la parte (b).
}
